
I. INTRODUCTION
Ultrasound is widely used because it is non-ionizing,
portable, and cost-effective. However, coherent wave interference introduces speckle noise that reduces contrast, obscures
lesion margins, and limits both human and automated interpretation. Speckle is not simple additive white noise; it is signal
dependent and spatially structured, which makes restoration
difficult without loss of clinically relevant texture.
Classical denoisers (Gaussian smoothing, wavelet thresholding, BM3D-like patch priors) can suppress noise, but often
at the cost of edge attenuation. Deep models, especially UNet architectures, can learn richer restoration priors from
data and preserve structure better than handcrafted filtering
in many image restoration tasks. Yet, most denoisers used
in practice remain point estimators: they output one image
without reporting how certain the model is about any pixel.
For medical deployment, this is a critical gap. If a model
is uncertain near lesion boundaries, the system should expose
this uncertainty rather than hide it. We therefore investigate
a denoising framework that pairs a compact U-Net with MC
dropout at inference to produce both a restored image and an
uncertainty map.
From a research methodology perspective, denoising quality
in medical imaging must be assessed under both fidelity and
reliability criteria. Fidelity concerns whether the output is close
to an underlying reference; reliability concerns whether the
model can indicate where this fidelity may degrade. Many
studies report only fidelity metrics and therefore do not address
the practical decision-support requirement in medical contexts.
In contrast, this paper explicitly treats uncertainty estimation
as a first-class research outcome.
Another challenge is evaluation realism. Reported denoising
gains can be inflated when experiments are restricted to one
corruption level or one split configuration. To reduce this risk,
we include robustness analysis under multiple train–test noise
pairings, provide transparent reporting about available versus
missing metrics in current artifacts, and frame conclusions
according to evidence strength rather than aspirational claims.
Research questions:
• RQ1: Does the proposed enhanced U-Net improve
restoration quality over a baseline U-Net model?
• RQ2: How robust is performance under train–test speckle
mismatch?
• RQ3: Do MC-dropout uncertainty maps localize difficult
restoration regions?
Contributions:
• A reproducible ultrasound denoising pipeline built around
a compact U-Net with optional wavelet-aware preprocessing and MC-dropout uncertainty.
• A structured evaluation protocol with aggregate results
and robustness sweeps across noise conditions.
• An uncertainty-aware analysis that complements fidelity
metrics and supports deployment-oriented interpretation.
• A detailed research reporting template that explicitly
documents limitations, threats to validity, and required
steps toward publication-grade evidence.
II. CLINICAL MOTIVATION AND PROBLEM STATEMENT
In breast ultrasound workflows, image quality directly impacts lesion visibility, contour readability, and confidence in
qualitative assessment. Speckle reduces conspicuity of subtle
boundaries and can induce texture confusion, especially in heterogeneous tissue regions. Any enhancement system intended
for clinical decision support should therefore satisfy at least
three constraints: preserve diagnostically relevant structures,
reduce corruption artifacts, and provide interpretable confidence indicators.
Let x ∈ [0, 1]H×W denote an idealized clean ultrasound
image and x˜ its observed noisy version. We seek a restoration
function fθ that maps noisy input to a denoised output while
preserving lesion boundaries and clinically relevant structure.
In parallel, we seek an uncertainty function uθ that assigns
higher values to visually ambiguous or difficult regions.
This perspective motivates a dual-output system: one output
for restoration, one output for reliability characterization. In
this work, MC dropout is used as an approximate Bayesian
mechanism to derive uθ through predictive variance.
III. RELATED WORK
Classical image denoising includes local smoothing, wavelet
shrinkage, and non-local patch methods. BM3D remains a
strong classical baseline in natural-image denoising due to
collaborative filtering in a transformed patch domain [4]. In
ultrasound, these methods reduce speckle but can blur lesion
boundaries when noise suppression is aggressive.
Deep denoisers such as CNN-based residual learning and UNet-style encoder–decoder networks model nonlinear image
priors and often outperform fixed-filter approaches [8], [1].
For medical images, skip connections are particularly valuable
because they preserve high-resolution structural detail during
reconstruction.
Uncertainty estimation in deep learning has gained importance for safety-critical settings. MC dropout approximates
Bayesian inference by keeping dropout active at test time and
sampling stochastic predictions [2]. In medical imaging, uncertainty maps are increasingly used to identify low-confidence
predictions and improve interpretability [9].
Relative to prior work, this paper focuses on an engineeringresearch bridge: an uncertainty-aware denoising model that is
both quantitatively evaluated and integrated into a deployable
backend/frontend stack for interactive analysis.
Recent restoration literature also highlights the role of
self-supervised and noise-model-aware training, including
Noise2Noise and related paradigms, for settings where clean
references are limited [10]. While our present implementation
uses supervised synthetic corruption, the broader methodological implication is that ultrasound enhancement may benefit
from hybrid supervision strategies in future iterations.
In medical AI practice, reproducibility and reporting transparency are increasingly recognized as essential for model
credibility. This includes explicit split definitions, hyperparameter disclosure, environment specification, and realistic
limitations analysis. The present paper adopts this reporting
stance and treats it as part of the research contribution rather
than supplementary material.
IV. PROPOSED METHODOLOGY
This section presents the complete proposed method as a
single end-to-end framework for ultrasound enhancement with
uncertainty estimation. The pipeline is designed to improve
image quality while simultaneously quantifying prediction
reliability.
A. Method Design Objectives
The proposed methodology is built around three design
objectives:
• Fidelity objective: suppress speckle noise while preserving lesion boundaries and structural detail.
• Reliability objective: provide pixel-wise uncertainty
cues to identify low-confidence regions.
• Deployability objective: support deterministic and
uncertainty-aware inference through a practical API+UI
workflow.
These objectives guide both the model design and the
evaluation protocol.
B. Core Method Components
The proposed method integrates four core components:
• Preprocessing module: grayscale conversion, resizing to
256 × 256, and intensity normalization to [0, 1].
• Optional wavelet preconditioning: configurable softthreshold denoising in the wavelet domain prior to deep
inference.
• U-Net restoration module: compact encoder-decoder
with skip connections and dropout-enabled convolution
blocks.
• Uncertainty module: MC-dropout inference that aggregates multiple stochastic predictions into mean output and
variance map.
The U-Net backbone follows a 1-32-64-128-256 channel
progression and uses mirrored decoder stages for structurepreserving reconstruction.
C. Stage-Wise Pipeline
The proposed methodology is organized into six stages:
1) Input Acquisition: A breast ultrasound image is provided by the user or evaluation dataset.
2) Preprocessing: The image is converted to grayscale (if
needed), resized to 256 × 256, and normalized to [0, 1].
3) Optional Wavelet Preconditioning: A configurable
Haar-wavelet soft-thresholding step is applied to reduce
high-frequency corruption before deep inference.
4) Deep Denoising: A compact U-Net performs image
restoration. Two operating modes are supported:
• Deterministic mode for low-latency enhancement.
• Stochastic MC-dropout mode for uncertainty-aware
enhancement.
5) Uncertainty Estimation: In MC mode, multiple
stochastic passes are aggregated to produce a mean denoised image and a pixel-wise uncertainty map derived
from prediction variability.
6) Output and Analysis: The system returns denoised
output, uncertainty map, and runtime metadata for interactive review in Streamlit and quantitative analysis
using PSNR/SSIM/MAE (where available).
D. Training and Inference Strategy
Training uses supervised noisy-clean pairs with configurable
objectives. The baseline uses mean squared error, while alternative losses (L1/frequency-aware) are available for exploratory studies. Checkpointing is performed in deploymentready and experiment-ready paths to support reproducibility.
At inference time, users can select between deterministic
and MC-dropout modes. Deterministic mode prioritizes latency, while MC mode prioritizes reliability-aware interpretation through prediction variability.
E. Why This Method Is Different
Compared with conventional deterministic denoising
pipelines, the proposed methodology explicitly couples
restoration and reliability estimation in one integrated
workflow. Instead of producing only a denoised image, the
method produces:
• a restored image for visual quality improvement,
• a pixel-wise uncertainty map for confidence-aware interpretation,
• runtime metadata for operational evaluation.
This dual-output design is central to the proposed contribution and directly supports the research questions on fidelity,
robustness, and uncertainty behavior.
F. Methodological Rationale
The proposed method is designed to address two practical
requirements that are often treated separately in ultrasound
enhancement studies:
• High-quality restoration: achieved through U-Net feature fusion and skip connections that preserve structural
detail.
• Reliability awareness: achieved through MC-dropout
uncertainty that indicates where model predictions are
less stable.
By combining both requirements in one deployable pipeline,
the method moves from purely image-enhancement performance toward clinically useful decision support.
G. Proposed Contribution Over Baseline Practice
Compared with a conventional deterministic denoiser, the
proposed methodology introduces:
• dual-output inference (restored image + uncertainty map),
• robustness-oriented evaluation across train–test noise mismatch,
• an implementation path that is directly deployable
through API endpoints and UI-based qualitative analysis.
TABLE I
EVALUATION TARGETS AND AVAILABILITY
Comparator Metric Availability Source
U-Net (MSE) PSNR, SSIM project table
U-Net + Hybrid + Noise PSNR, SSIM project table
MAE summary not exported UI/per-image only
This integrated design is the central methodological contribution of the paper.
V. EXPERIMENTAL DESIGN
A. Dataset and Splits
Experiments use BUSI-style ultrasound images [5]. The
training script performs a seeded random split with default
validation ratio 0.2. Images are resized to 256 × 256 for all
runs.
To maintain experimental consistency, all comparisons use
identical preprocessing and image geometry. The same core
architecture is reused across baseline and enhanced settings,
so observed gains are attributable to training/inference strategy
differences rather than arbitrary architectural inflation.
B. Protocol and Hypotheses
We evaluate three hypotheses aligned with RQ1–RQ3:
• H1: Enhanced U-Net achieves higher PSNR/SSIM (and
lower MAE where available) than baseline U-Net.
• H2: Performance decreases under train–test noise mismatch, but remains acceptable within moderate mismatch.
• H3: High uncertainty concentrates in visually challenging
regions (speckle-heavy background and boundary transitions).
C. Baselines and Comparators
We compare model variants currently available in the project
artifacts. Classical comparators are included for context in
the literature discussion, while the quantitative table reports
measured values from available logs.
The current artifact-backed comparison is intentionally conservative: only values with explicit exported evidence are
treated as quantitative claims. This avoids over-interpretation
and aligns with robust scientific reporting practice.
D. Metrics
We report PSNR, SSIM [6], and MAE when available.
PSNR summarizes reconstruction fidelity relative to mean
squared error, SSIM evaluates structural similarity, and MAE
provides an interpretable absolute-error measure.
E. Implementation Details
Default retraining parameters are: 12 epochs, batch size 8,
learning rate 10−3, Adam optimizer, and synthetic training
noise standard deviation 0.03. Inference supports FP32 and
optional FP16. MC sample count is user-controlled in the
interface (default 10).
Training checkpoints are persisted in two locations to support both deployment and experiment reproducibility. Runtime
behavior (including per-request latency and model-loading
status) is exposed by the backend and surfaced in the user
interface for transparent operational profiling.
F. Reproducibility Protocol
The implementation includes deterministic seeding for
Python, NumPy, and PyTorch random number generators. This
supports repeatable split generation and improves comparability across retraining runs. A publication-grade extension
should additionally report exact package versions, hardware
details, and confidence intervals across multiple random seeds.
VI. SYSTEM ARCHITECTURE
The system is implemented as a full-stack application with
three main components. The frontend is built with Streamlit
and provides an interactive interface for image upload, visualization, and result inspection. The backend is implemented
using Flask and exposes inference endpoints for deterministic
denoising and uncertainty-aware prediction. The model itself
is implemented in PyTorch.
The processing flow is: input image → preprocessing →
U-Net inference → optional Monte Carlo dropout sampling
→ denoised output + uncertainty map.
Fig. 1. System flow diagram for ultrasound denoising with uncertainty
estimation.
TABLE II
LAYER-WISE ARCHITECTURE OF THE PROPOSED U-NET
Stage Operation Output Tensor Size
Input Grayscale image 1 × 256 × 256
Enc-1 DoubleConv (3 × 3), Dropout 32 × 256 × 256
Pool-1 MaxPool (2 × 2) 32 × 128 × 128
Enc-2 DoubleConv (3 × 3), Dropout 64 × 128 × 128
Pool-2 MaxPool (2 × 2) 64 × 64 × 64
Enc-3 DoubleConv (3 × 3), Dropout 128 × 64 × 64
Pool-3 MaxPool (2 × 2) 128 × 32 × 32
Bottleneck DoubleConv (3 × 3), Dropout 256 × 32 × 32
Up-3 Transposed Conv + concat(Enc-3) 256 × 64 × 64
Dec-3 DoubleConv (3 × 3), Dropout 128 × 64 × 64
Up-2 Transposed Conv + concat(Enc-2) 128 × 128 × 128
Dec-2 DoubleConv (3 × 3), Dropout 64 × 128 × 128
Up-1 Transposed Conv + concat(Enc-1) 64 × 256 × 256
Dec-1 DoubleConv (3 × 3), Dropout 32 × 256 × 256
Output Conv (1 × 1) + Sigmoid 1 × 256 × 256
Fig. 2. System architecture diagram showing module interactions and UI
integration.
A. Proposed Network Architecture (Layer-Wise)
To make the model definition fully reproducible, Table II
provides the layer-wise architecture used in this work. Input
resolution is fixed at 256 × 256 with one grayscale channel.
The architecture follows a 3-level encoder, bottleneck, and
symmetric decoder with skip concatenations.
The final sigmoid layer constrains output intensities to [0, 1],
matching the normalized target domain. Dropout is retained
in all convolutional blocks so the same architecture supports
both deterministic inference and MC-dropout sampling for
uncertainty estimation.
Architecturally, the separation between API-layer inference
and UI-layer analysis supports research operations. The back-
end can be benchmarked independently, while the frontend
can be used for qualitative review, uncertainty inspection, and
report generation. This modularity reduces coupling between
modeling code and presentation logic, which is important for
experimental maintainability.
VII. RESULTS
A. Quantitative Comparison (RQ1)
Table IV reports aggregate metrics from project artifacts.
The enhanced model improves PSNR by +6.98 dB and SSIM
by +0.254 over the baseline U-Net (MSE), supporting H1.
These gains are substantial for ultrasound denoising and indicate stronger restoration of both intensity fidelity and structural
consistency.
The relative improvement is also meaningful when interpreted proportionally. The SSIM increase from 0.488 to 0.742
corresponds to approximately 52% relative gain over baseline
structural similarity. In restoration studies, improvements of
this magnitude generally indicate not only reduced noise but
improved preservation of image organization.
Because MAE was not exported in the same aggregate table,
we do not over-claim MAE improvements at dataset level. This
is an explicit reporting limitation rather than an omitted result.
B. Robustness Under Noise Mismatch (RQ2)
Table V and Fig. 5 summarize train–test mismatch behavior.
Best performance occurs at low matched noise (0.2/0.2), with
a monotonic decline as test noise increases. At matched 0.4/0.4
and 0.6/0.6 settings, performance remains comparatively stable, indicating reasonable adaptation to training corruption
level. These observations support H2: mismatch hurts performance, but degradation is gradual rather than catastrophic.
C. Uncertainty Behavior (RQ3)
The MC-dropout output exposes spatially varying uncertainty rather than a single scalar confidence. In qualitative
examples, high-variance regions co-localize with noisy texture
and boundary transitions, while homogeneous regions show
lower variance. This behavior is consistent with H3 and
suggests practical value for decision support: clinicians can
visually inspect where restoration is less trustworthy.
A useful practical interpretation is uncertainty-guided attention: instead of trusting the denoised image uniformly, downstream users can inspect high-variance areas first. This does
not replace diagnosis but can improve review prioritization.
D. Qualitative Findings
Fig. 3 shows clear reduction of speckle granularity with
preserved large-scale structure. Fig. 4 highlights residual error
concentration around challenging structures. Together, these
visual findings align with the quantitative PSNR/SSIM gains.
Additional qualitative snapshots from project notebook exports are provided in Fig. 6 and Fig. 7. These examples were
generated by the same project pipeline and provide supplementary visual evidence of consistent enhancement behavior.
TABLE III
CORE TRAINING AND INFERENCE CONFIGURATION
Component Setting
Input size 256 × 256
Input channels 1 (grayscale)
U-Net dropout 0.1
Optimizer Adam
Initial learning rate 10−3
Batch size 8
Epochs 12
Validation split 0.2
Training noise std. 0.03
MC dropout samples 2–30 in UI, 10 default
Inference precision FP32 / optional FP16
TABLE IV
SUMMARY OF RECORDED RESULTS
Method PSNR (dB) SSIM MAE
U-Net (MSE) 20.04 0.488 Not exported
U-Net + Hybrid + Noise 27.02 0.742 Not exported
E. Statistical Rigor and Current Gaps
The current repository results represent a strong pilot benchmark but do not yet include repeated-run confidence intervals,
fold-level variance, or significance testing. Therefore, conclusions are framed as evidence-supported but preliminary. A
full statistical protocol (multi-seed training, bootstrap confidence intervals, paired tests per image) is recommended for
publication-grade claims.
Specifically, a future statistical pipeline should include: (i)
at least 3–5 independent seeds per setting, (ii) per-image metric distributions, (iii) paired hypothesis tests between model
variants, and (iv) uncertainty-error association analysis (e.g.,
correlation between predictive variance and absolute residual
magnitude).
VIII. DISCUSSION
The experimental evidence supports all three research questions at pilot-study level. For RQ1, the improved PSNR/SSIM
over baseline indicate clear restoration gains. For RQ2, robustness curves confirm expected sensitivity to mismatch while
preserving usable quality across tested conditions. For RQ3,
uncertainty maps add clinically relevant context by flagging
potentially unreliable regions.
Importantly, this work emphasizes evaluation discipline.
We report only metrics that are available in project artifacts,
explicitly identify missing aggregate MAE exports, and avoid
unsupported statistical claims. This reporting posture is crucial
for credible medical-AI research.
From a systems perspective, the architecture is deployable:
a Flask inference API and Streamlit analysis interface already
exist. This allows direct translation from offline experiments
to interactive validation workflows.
The key research implication is that restoration quality and
reliability should be evaluated jointly. A model with high
PSNR but no uncertainty channel may still be unsafe in
ambiguous cases. Conversely, uncertainty estimates without
strong restoration quality are also insufficient. The present
framework demonstrates a practical compromise: strong denoising gains plus an interpretable reliability surrogate.
A second implication concerns methodology: transparent reporting of missing metrics and unresolved statistical questions
should be treated as a strength, not a weakness, in early-stage
medical AI research. It prevents overstatement and clarifies
what evidence is currently established versus what remains to
be validated.
IX. DEPLOYMENT CONSIDERATIONS
From a deployment perspective, the current system is practical because it already exposes a Flask inference layer and
a Streamlit interface. This means the method is not only a
training-time model but a working application that can be
extended for clinical experimentation, batch processing, or
web-based review.
At the same time, the reported results should be interpreted
as project-level experimental evidence rather than a final
clinical benchmark. Additional validation on larger datasets
and external validation splits would be necessary before any
clinical claims.
In real-world deployment, uncertainty output could be used
to trigger quality-control pathways, such as flagging lowconfidence frames for manual review or requesting additional
imaging views. This integration pathway is one of the strongest
practical benefits of uncertainty-aware restoration compared to
deterministic enhancement alone.
Operationally, model monitoring should include drift detection, periodic recalibration checks, and failure-case logging. Ultrasound domain shift can arise from scanner vendor
differences, acquisition protocols, and operator variability;
therefore, post-deployment surveillance is essential.
X. LIMITATIONS AND THREATS TO VALIDITY
Internal validity: The present benchmark does not yet
include multi-seed variance estimates or repeated cross-
validation. Observed improvements may contain run-to-run
sensitivity.
Construct validity: Speckle is simulated with a multiplicative Gaussian model. Real ultrasound artifacts can deviate from
this assumption, potentially changing relative performance.
External validity: Results are based on BUSI-style data
and internal project artifacts. Generalization to other organs,
scanners, and acquisition protocols remains to be tested.
../results/images/before_after.png
Fig. 3. Representative before-and-after denoising output from the project
results.
../results/images/diff_map.png
Fig. 4. Difference or uncertainty-related visualization generated by the project
pipeline.
validation. Measurement validity: MAE is available per-image in
tooling but not fully exported in current aggregate tables,
creating an incomplete tri-metric summary.
Comparator validity: Current quantitative reporting is
strongest for project-native model variants, while some classical comparators are literature contextual rather than fully
re-benchmarked under identical implementation conditions in
this artifact.
Visualization validity: Uncertainty maps are qualitatively
../experiments/notebook_run/exports/preview_20260Fig. 7. Additional qualitative preview generated from notebook-run export
informative, but quantitative calibration metrics (e.g., riskcoverage, expected calibration error in task-aligned formulations) are not yet reported.
XI. ETHICAL AND TRANSLATIONAL CONSIDERATIONS
Medical enhancement models should be framed as assistive
tools, not autonomous diagnostic systems. Denoising may improve readability but can also alter visual perception; therefore,
transparent indication of model uncertainty is important to
reduce over-reliance.
For translational use, governance concerns include traceability of model versions, auditability of inference outputs,
and documentation of known failure modes. Future extensions
should include human-factor studies to evaluate how radiologists interpret uncertainty overlays and whether these overlays
improve decision quality or only increase cognitive load.
XII. REPRODUCIBILITY AND OPEN SCIENCE STATEMENT
This paper is grounded in executable project artifacts including model code, API endpoints, and result tables. To
improve reproducibility for external readers, the next release
should include: (i) frozen dependency manifests, (ii) scripted
experiment runners for all reported tables, (iii) seed lists and
deterministic flags, and (iv) machine-readable metadata for
each result file.
We emphasize that reproducibility is not binary. The present
work is reproducible at prototype level and transparently
identifies what is required to reach archival, publication-grade
reproducibility.
XIII. CONCLUSION
This study demonstrates that uncertainty-aware U-Net denoising is a credible approach for ultrasound enhancement
under speckle corruption. The enhanced model shows strong
gains over baseline in available aggregate metrics, and MC
dropout provides a useful spatial confidence signal that complements deterministic restoration. The work contributes both
a measurable denoising improvement and a deployable analysis pipeline.
More broadly, the study argues for a dual-objective perspective in medical image restoration: optimize image fidelity
while quantifying predictive reliability. This perspective is
essential for trustworthy AI in clinical imaging and should
become standard in future denoising benchmarks.
XIV. FUTURE WORK
Future work will focus on publication-grade validation: (1)
multi-seed and fold-wise statistical testing, (2) complete MAE
export in all summary tables, (3) external dataset evaluation,
(4) comparison against stronger classical/deep baselines under
identical protocols, and (5) uncertainty calibration studies
(e.g., uncertainty-error correlation and coverage metrics).
Additional technical directions include domain-adaptive
training for scanner variation, uncertainty-aware loss shaping,
and integration with segmentation/classification pipelines to
evaluate whether denoising-plus-uncertainty improves downstream clinical tasks.
ACKNOWLEDGMENT
The authors would like to acknowledge the ultrasound
image dataset and the project implementation artifacts used
to prepare this manuscript draft.
REFERENCES
[1] O. Ronneberger, P. Fischer, and T. Brox, “U-Net: Convolutional Networks for Biomedical Image Segmentation,” in Medical Image Computing and Computer-Assisted Intervention, 2015.
[2] Y. Gal and Z. Ghahramani, “Dropout as a Bayesian Approximation:
Representing Model Uncertainty in Deep Learning,” in ICML, 2016.
[3] I. Goodfellow, Y. Bengio, and A. Courville, Deep Learning. MIT Press,
2016.
[4] K. Dabov, A. Foi, V. Katkovnik, and K. Egiazarian, “Image Denoising
by Sparse 3-D Transform-Domain Collaborative Filtering,” IEEE Transactions on Image Processing, 2007.
[5] W. Al-Dhabyani, M. Gomaa, H. Khaled, and A. Fahmy, “Dataset of
Breast Ultrasound Images,” Data in Brief, 2020.
[6] Z. Wang, A. C. Bovik, H. R. Sheikh, and E. P. Simoncelli, “Image
quality assessment: from error visibility to structural similarity,” IEEE
Transactions on Image Processing, 2004.
[7] N. Jayant and P. Noll, Digital Coding of Waveforms. Prentice Hall, 1984.
[8] K. Zhang, W. Zuo, Y. Chen, D. Meng, and L. Zhang, “Beyond a Gaussian Denoiser: Residual Learning of Deep CNN for Image Denoising,”
IEEE Transactions on Image Processing, 2017.
[9] A. Kendall and Y. Gal, “What Uncertainties Do We Need in Bayesian
Deep Learning for Computer Vision?,” in NeurIPS, 2017.
[10] J. Lehtinen et al., “Noise2Noise: Learning Image Restoration without
Clean Data,” in ICML, 2018.
[11] A. Buades, B. Coll, and J.-M. Morel, “A Non-Local Algorithm for Image
Denoising,” in CVPR, 2005.
[12] X. Mao, C. Shen, and Y.-B. Yang, “Image Restoration Using Very
Deep Convolutional Encoder-Decoder Networks with Symmetric Skip
Connections,” in NeurIPS, 2016.
[13] K. He, X. Zhang, S. Ren, and J. Sun, “Deep Residual Learning for
Image Recognition,” in CVPR, 2016.
[14] P. Isola, J.-Y. Zhu, T. Zhou, and A. A. Efros, “Image-to-Image Translation with Conditional Adversarial Networks,” in CVPR, 2017.
[15] G. Huang, Z. Liu, L. van der Maaten, and K. Q. Weinberger, “Densely
Connected Convolutional Networks,” in CVPR, 2017.
